\section{Legacy SDL GPU Chapter Material}
\label{ch:sdlgpu-backend}

\texttt{SDL\_GPU} is a CNA graphics renderer, and one this book itself was slow to
give a proper account of --- Chapter~\ref{ch:roadmap} names that delay honestly as a gap in
this book, not in CNA's own code. This chapter closes it: a real, Linux/Vulkan-driver-verified
renderer implementing the full \cnaclass{IGraphicsRenderer} contract
(Chapter~\ref{ch:backend-architecture}), built on top of SDL3's own \texttt{SDL\_gpu} API
rather than a native vendor API of its own.

\section{Why this renderer costs nothing new}

\texttt{SDL\_gpu} dispatches internally to Vulkan, D3D12, or Metal depending on platform and
driver availability, and CNA's own \texttt{plan\_sdlgpu.md} states plainly why this renderer was
worth adding on top of the hardware-backed renderers the project already had: SDL3 is
already a vendored dependency of this project (\texttt{third\_party/SDL}), and
\texttt{SDL\_gpu.c} is already compiled into every prebuilt SDL3 package this repository
produces --- unlike Vulkan (the Vulkan SDK/loader) or WebGPU (\texttt{wgpu-native}, a separate
library), standing this renderer up required zero new third-party dependency, only new CNA-side
code against an API already sitting in the build. The main class is
\texttt{CNA::\allowbreak{}Internal::\allowbreak{}Renderers::\allowbreak{}SdlGpu::\allowbreak{}SdlGpuRenderer}
\index{CNA::Internal::Renderers::SdlGpu::SdlGpuRenderer@CNA::Internal::Renderers::\allowbreak{}SdlGpu::SdlGpuRenderer}.
Its source lives under
\texttt{modules/renderers/sdl-gpu/src/}. Build the target
\texttt{cna\_\allowbreak{}renderer\_\allowbreak{}sdl\_\allowbreak{}gpu}. Select this renderer with
\texttt{-DCNA\_\allowbreak{}GRAPHICS\_\allowbreak{}RENDERER=SDL\_\allowbreak{}GPU}.
The project's own naming table is explicit that this class name is deliberately distinct from
the pre-existing \cnaclass{SdlRenderer} in \texttt{SdlRenderer/} --- the older
\texttt{SDL\_Renderer} 2D API and the newer \texttt{SDL\_gpu} API are unrelated, and this project
was careful not to let the similar SDL naming collide.

\section{Precompiled bytecode only, in a fixed resource-binding order}

\texttt{SDL\_gpu} does not accept GLSL or HLSL source text at runtime the way
\cnaclass{ShaderEffect}'s own public contract promises a game --- \texttt{SDL\_CreateGPUShader}
takes only \texttt{SDL\_GPUShaderFormat}-tagged precompiled bytecode: SPIR-V for the Vulkan
driver, DXBC/DXIL for the D3D12 driver, MSL/\texttt{metallib} for the Metal driver. A further
constraint worth knowing before assuming any existing shader is reusable as-is: every stage's
samplers, storage textures, storage buffers, and uniform buffers must be declared in a fixed,
HLSL-style binding order (\texttt{t[n]} / \texttt{s[n]} in one address space, \texttt{u[n]} in
another, \texttt{b[n]} in a third) --- a different convention from the raw
\texttt{layout(set=, binding=)} numbering the existing \cnaclass{VulkanRenderer}'s own
GLSL shaders use. This means Vulkan's compiled SPIR-V is a useful \emph{algorithmic} reference
for this renderer's own shader authoring (the same lighting math, alpha-test logic, and
dual-texture blend) but not directly reusable output --- every shader's binding layout had to
be re-authored specifically for \texttt{SDL\_gpu}'s own convention, confirmed directly against
the real \texttt{.glsl} sources under this renderer's own
\texttt{modules/renderers/sdl-gpu/src/shaders/}
directory (23 files today, one pair per stock effect --- \texttt{colored3d},
\texttt{alpha\_test3d}, \texttt{dual\_texture3d}, \texttt{env\_map3d}, \texttt{lit\_textured3d},
\texttt{pbr3d} / \texttt{pbr\_skinned3d}, \texttt{skinned3d} / \texttt{skinned\_colored3d}, and more).

\subsection{A worked example: the Y-flip bug a real screenshot caught}

This renderer's own first real milestone --- a textured, tinted, rotated, multi-flip sprite
scene surviving a resize --- is a concrete instance of a theme this book returns to
repeatedly (Chapter~\ref{ch:testing-philosophy}): a ``did it throw'' test alone cannot catch a
wrong-but-plausible image. \texttt{SDL\_gpu}'s Vulkan driver flips clip-space Y internally, for
cross-renderer NDC consistency reasons of its own --- the vertex shader's first, naively-ported
NDC computation (\texttt{gl\_Position = vec4(ndc, 0, 1)}, the same shape every other renderer's
own sprite vertex shader already uses) compiled cleanly, ran without error, and rendered every
single sprite upside down, a texture's top row appearing at the bottom of the screen. Only a
real captured screenshot caught it; the fix negates Y explicitly
(\texttt{vec4(ndc.x, -ndc.y, 0, 1)}), confirmed against an isolated single-sprite diagnostic
tool kept permanently in the tree (\texttt{modules/renderers/sdl-gpu/examples/\allowbreak{}sdlgpu\_diag\_single\_sprite.cpp}, the same
``keep the diagnostic, don't throw it away once the bug is fixed'' precedent
Chapter~\ref{ch:direct3d-backends} already covers for \texttt{cna\_diag\_d3d12\_swapchain}).

\section{SkinnedEffect's storage-buffer workaround: a real, undocumented push-uniform cap}

Wiring \cnaclass{SkinnedEffect} through this renderer surfaced a genuine hardware/driver
constraint rather than a CNA design choice: \texttt{SDL\_PushGPUVertexUniformData}, the
mechanism every other stock effect above uses to deliver its per-draw uniforms, has a real,
undocumented cap of roughly 4096 bytes per slot on this renderer's own Vulkan driver --- found by
a real empirical binary-search spike, not read out of any SDL3 header or release note, since
the limit is not documented anywhere in \texttt{SDL\_gpu.h} itself. A 72-bone skin palette
(\(72 \times 64\) bytes \(= 4608\) bytes) exceeds that cap outright, so
\cnaclass{SkinnedEffect}'s bone matrices are uploaded through a real GPU storage buffer
(\texttt{SDL\_BindGPUVertexStorageBuffers}) instead of a uniform push --- every other stock
effect's smaller per-draw payloads (world/view/projection, \texttt{DiffuseColor}, alpha-test
parameters, light parameters) stay on the ordinary uniform-push path, since none of them come
close to the cap. This is exactly the kind of real, hardware-shaped constraint this book asks
every renderer chapter to name precisely rather than paper over with ``uniforms are uploaded
per draw'' as if the mechanism were uniform across every payload size.

\section{Custom ShaderEffect: a real runtime compile, built without the dev package installed}

Custom \cnaclass{ShaderEffect} support (Chapter~\ref{ch:shader-fx-gap}) is real on this renderer
too, and its own implementation is worth knowing about for a reason beyond the feature itself.
Since \texttt{SDL\_gpu} accepts only precompiled bytecode, \cnaclass{SdlGpuEffectRenderer} bridges
\cnaclass{ShaderEffect}'s public GLSL-source contract to that requirement with a genuine runtime
\texttt{libshaderc} GLSL-to-SPIR-V compile, linked directly into the renderer binary rather than
merely invoked at build time. Reading \texttt{SdlGpuRenderer.cpp} directly shows a small,
telling detail: this dev environment has no \texttt{libshaderc-dev} package installed, so there
is no real \texttt{shaderc.h} to include --- the renderer's own source hand-declares the exact
handful of \texttt{shaderc\_*} C entry points it needs. These include
\texttt{shaderc\_\allowbreak{}compiler\_\allowbreak{}initialize},
\texttt{shaderc\_\allowbreak{}compile\_\allowbreak{}into\_\allowbreak{}spv}, and
\texttt{shaderc\_\allowbreak{}result\_\allowbreak{}get\_\allowbreak{}bytes}. Their
\texttt{extern "C"} function signatures, shader-kind values, and optimization-level values are
hand-authored as named constants rather than pulled from a missing header.
This is a real, working engineering accommodation to a genuinely constrained dev environment,
not a shortcut around the feature itself --- the compiled shader still goes through the real
linked \texttt{libshaderc.so}, only the header describing its C ABI is hand-authored instead of
vendored. Building this custom-shader path also surfaced a genuine finding one layer over: the
existing \cnaclass{VulkanRenderer}'s own \cnaclass{ShaderEffect} support does not
actually runtime-compile GLSL text at all, despite \cnaclass{ShaderEffect}'s own documented
contract implying every renderer does --- a real, previously-unflagged discrepancy between one
renderer's implementation and the class's own stated cross-renderer promise, left for
Chapter~\ref{ch:vulkan-backend} to account for on its own terms rather than resolved here.

\subsection{The \texttt{.cnj} fixture exposes the remaining shader-compatibility boundary}

The runtime compiler does not make every GLSL text accepted elsewhere in CNA portable to this
renderer. The real \texttt{CnjEffectTest.LoadsRealCnjFixture} content test writes an ordinary
\texttt{.cnj} effect manifest plus a GLSL ES~3.00 vertex/fragment pair: \texttt{\#version 300 es},
unqualified stage outputs/inputs, and loose \texttt{uniform mat4 projection} /
\texttt{uniform sampler2D texture1} declarations. Under \texttt{SDL\_GPU}, that fixture fails
before a pixel is drawn. Its source is not in the SPIR-V-compatible dialect required by the
runtime compiler and by the resource counts CNA supplies to \texttt{SDL\_CreateGPUShader}.

The passing \texttt{SdlGpu\_ShaderEffect} regression makes the required contrast concrete. Its
two sources use desktop \texttt{\#version 450}, explicit \texttt{layout(location=...)} qualifiers
for every vertex attribute and stage I/O, and the exact fixed binding spaces this renderer
declares: the fragment sampler is \texttt{layout(set=2, binding=0)}, while the per-draw values
live in a named uniform block at \texttt{set=1} for the vertex stage and \texttt{set=3} for the
fragment stage. It compiles both stages at runtime, draws a white texture into a
\cnaclass{RenderTarget2D}, and reads back the effect's own blue-tinted value; repeating the
same draw without the effect must read white. Thus the supported dialect and binding contract
are proven by pixels, not merely by \texttt{IsEffectValid()}.

This is deliberately an open compatibility gap, not evidence that \cnaclass{ShaderEffect} is a
no-op on SDL\_GPU. Closing it needs a separately designed decision: either migrate the shared
\texttt{.cnj} fixture/source convention to the stricter portable dialect, or teach this renderer
to translate its existing loose GLSL into the blocks, locations, and bindings SDL\_gpu requires.
Neither is a safe one-line relaxation; the first changes content expectations across renderers,
and the second is a real shader-translation layer. The chapter therefore records the narrower,
honest current contract: custom effects work when authored for SDL\_GPU's binding dialect, while
the generic \texttt{.cnj} fixture is not yet cross-renderer portable to it.

\section{Pipeline state is snapshotted at draw time, except where it is deliberately not}

SDL\_gpu graphics pipelines are immutable enough that CNA cannot simply bind an XNA state object
at the end of a deferred frame and hope it represents every earlier draw. Each queued sprite or
3D draw therefore captures a \texttt{RenderStateSnapshot}: blend factors/functions, cull and
fill mode, stencil operations/masks/reference, and the ordinary depth-test fields travel with
the command that was issued under them. Pipeline lookup folds the pipeline-relevant portion into
a \texttt{size\_t} hash together with primitive topology, target color format, and sample count.
This replaced an earlier hand-packed integer key, whose fixed bit budget would have become
silently collision-prone as state dimensions were added.

The detail that keeps the cache both correct and finite is conditional hashing. Disabled blend
does not hash its six blend factors/functions; disabled stencil does not hash operations or masks;
and disabled two-sided stencil does not hash the counter-clockwise fields. Two commands that
produce the same pipeline consequently share it even if their irrelevant XNA state-object fields
differ, while a meaningful state change cannot reuse a pipeline with stale SDL\_gpu descriptors.

\begin{sourcenote}
The target format included in that key is not the public backbuffer-format field. Each frame
queries SDL for the swapchain texture's actual format and builds the matching pipeline.
Independently, construction chooses D24S8 or D32S8 by a real device-support query and later allocates that
combined attachment whenever a backbuffer frame needs it. A public
\texttt{DepthFormat::None} does not suppress it; absence means the device supported neither
combined format. Fullscreen remains an SDL-window operation outside this renderer hook, and
swapchain acquisition follows the resulting window without selecting formats from the public
enums.
\end{sourcenote}
Stencil reference itself is intentionally \emph{not} pipeline state: the captured value is sent
per draw through \texttt{SDL\_SetGPUStencilReference()}, the missing call that an adversarial
pixel test once exposed. A left-half stencil write followed by an Equal-reference reveal now
proves that this dynamic value reaches the GPU rather than merely being stored in CNA.

Viewport, scissor, and blend constants follow the same per-command rule. Every queued draw
reference snapshots the active viewport rectangle/depth range, scissor rectangle and enable bit,
and blend factor; replay applies them through \texttt{SDL\_SetGPUViewport},
\texttt{SDL\_SetGPUScissor}, and \texttt{SDL\_SetGPUBlendConstants} immediately before that
draw. Enabled scissors are clamped to the target; disabled or degenerate ones expand to its full
extent. The shared deferred-viewport/scissor suites queue several changes without a flush and
cover backbuffer, 2D-target, and cube-face routes, so a final live state cannot satisfy them by
accident.

Depth bias is pipeline-static in SDL\_gpu 3.5 rather than dynamic. CNA normalizes both
\cnaclass{RasterizerState.DepthBias} and \texttt{SlopeScaleDepthBias} by primitive topology,
includes their bits in every affected pipeline key, and fills SDL's constant and slope factors
when creating the pipeline. The dedicated test combines structural key checks, differential
depth pixels, and validation-fatal execution. Samplers likewise carry the full filter ordinal,
U/V address modes, and clamped \texttt{MaxAnisotropy} through both cache key and native
descriptor; anisotropy is no longer merely a stored capability bit.

\section{Deferred rendering still preserves the game's draw order}

Deferring all SDL\_gpu work until \texttt{Present()} once introduced a correctness bug hidden by
the renderer's family-specific command vectors: it rendered every 3D family first and every sprite
last, regardless of the order the game issued them. A background sprite, 3D model, and HUD sprite
therefore became background/model/HUD only by accident; even the opposite sprite--then--3D order
was replayed as 3D--sprite. This is visible whenever later opaque or blended content should cover
earlier content, not merely a performance detail.

Every \texttt{Queue*Draw()} and \texttt{QueueSprite()} call now appends a compact kind/index
reference beside its family-specific command. At render time,
\texttt{RenderQueuedDraws()} makes one chronological pass over those references and dispatches
the existing per-family issue routine only when its command belongs to the current target. No
sort is needed: append order is already the public draw-call order. Pipeline rebind avoidance is
then tracked across all kinds rather than only sprites, so fixing correctness does not require
forgetting the ordinary bind cache.

\texttt{SdlGpu\_DrawOrder} proves both directions with no alpha arithmetic ambiguity. It renders
an opaque red full-target sprite and an opaque green full-target 3D quad into separate targets
with depth disabled. Sprite then 3D must read back green; 3D then sprite must read back red.
The former fixed-family implementation produced red in both cases because sprites always ran
last, so the pair is a discriminating regression rather than a no-crash smoke check.

\section{Present timing belongs to the device manager, not a test-only renderer poke}

The SDL\_GPU tests establish a small but important rule about CNA's initialization lifecycle.
\cnaclass{GraphicsDeviceManager} defaults \texttt{SynchronizeWithVerticalRetrace} to true,
which is the right XNA-compatible default for an ordinary game. On this project's virtual test
display, however, there is no useful vertical-retrace signal; leaving it enabled made each frame
wait roughly one second. A 60- or 120-frame GPU regression test therefore consumed its CTest
budget without doing more useful rendering. The correct setup, used by all 22 current
SDL\_GPU examples and diagnostics, is public configuration made before
\texttt{Game::DoInitialize()} creates or resets the device:

\begin{lstlisting}
gdm_ = std::make_unique<GraphicsDeviceManager>(this);
gdm_->setPreferredBackBufferWidthProperty(64);
gdm_->setPreferredBackBufferHeightProperty(64);
gdm_->setSynchronizeWithVerticalRetraceProperty(false);
\end{lstlisting}

This order matters. Earlier tests called \texttt{SdlGpuRenderer::SetSwapInterval(0)}
directly as a local workaround. During initialization, though, the device manager converts its
still-true public property to \texttt{PresentInterval::One}; \cnaclass{GraphicsDevice::Reset}
then forwards that presentation interval to every renderer. The correct global forwarding was
itself a real repair: before it landed, the reset path did not call
\texttt{IGraphicsRenderer::SetSwapInterval()} at all. Once repaired, it rightly overwrote the
private pre-initialization poke, revealing the test timeout instead of preserving an accidental
ordering dependency.

For this renderer, interval zero asks SDL\_gpu for \texttt{SDL\_GPU\_PRESENTMODE\_IMMEDIATE},
falling back to \texttt{MAILBOX} if Immediate is unavailable; any positive interval selects
VSync. SDL\_gpu has no half-rate present mode, so XNA's \texttt{PresentInterval::Two} has the
same VSync outcome as \texttt{One}. That makes the example's explicit false a test-environment
choice, not a claim that games should globally disable synchronization: it uses the same public
property a game would use, at the phase of the lifecycle where CNA will not subsequently replace
it.

\section{Validation mode turned two tolerated violations into real fixes}

This renderer now passes SDL\_gpu's \texttt{debug\_mode} flag from CNA's ordinary build mode:
debug builds request validation and release builds do not. That deliberately modest toggle paid
for itself immediately. With validation previously hardcoded off, two API-contract violations
had appeared to work; with it on, each could hang the Vulkan driver rather than merely print a
warning. The lesson is more useful than a generic ``enable validation'' recommendation: a
renderer's release-mode pixels can look plausible while resource descriptions are still invalid.

First, the MSAA implementation for \cnaclass{RenderTargetCube} used a six-layer 2D-array
texture. SDL\_gpu forbids multisampling on an array texture. CNA now uses one single-layer,
multisampled 2D attachment for whichever cube face is currently rendering and resolves it into
the actual cube texture; cycling is permitted only for that temporary attachment, never for the
cube texture that must retain all six faces. Second, the initial automatic-mipmap implementation
for plain \cnaclass{Texture3D} and \cnaclass{TextureCube} requested only \texttt{SAMPLER} usage,
although SDL\_gpu's generator requires \texttt{COLOR\_TARGET} too. Widening the usage exposed a
third, deeper failure: SDL's Vulkan path created invalid 2D views and blits for depth planes that
no longer exist in smaller 3D mip levels.

That third finding is no longer open. \texttt{REMED-GFX-099} removed the invented automatic
generation behavior from plain \cnaclass{Texture3D}: XNA/FNA's contract allocates an authored
chain and lets each \texttt{SetData(level,...)} call populate the named level. The resource is
therefore \texttt{SAMPLER}-only again, with no target views or generator blits to violate the
Vulkan contract. A validation-fatal regression now covers exact level-zero and nonzero-mip
round trips, partial volumes, distinguishable depth planes, forward/reverse write order, NPOT
dimensions, and two alternating resources. Plain cube textures retain their separately tested
per-face behavior. This sequence is why validation must be treated as evidence with scope: it
first exposed tolerated violations, then forced the design back to the actual public resource
contract rather than merely suppressing their diagnostics.

\section{Render targets: a real use-after-free, closed before it ever shipped}

\cnaclass{RenderTarget2D} and \cnaclass{RenderTargetCube} support (with real multisampling and
mip regeneration) surfaced a genuine lifetime hazard while its own test was still being written,
not after: this renderer defers its actual render pass to \texttt{Present()} time, so a render
target destroyed \emph{before} that deferred pass executes was a real use-after-free the moment
anything still referenced its underlying GPU state. The fix,
\cnaclass{SdlGpuRenderTarget2DState} / \cnaclass{SdlGpuRenderTargetCubeState}, is a pair of
\texttt{shared\_ptr}-owned structs holding the actual GPU state independently of the
\cnaclass{SdlGpuRenderTargetRenderer} / \cnaclass{SdlGpuRenderTargetCubeRenderer} wrapper's own
C\texttt{++} object lifetime --- a short-lived, local render target destroyed mid-\texttt{Draw()}
now keeps its pending content alive long enough to render correctly and remain safely sampleable
elsewhere the same frame. The bug was confirmed via a real segfault in
\texttt{sdlgpu\_mrt\_test.cpp} before the fix landed, not merely reasoned about in the abstract.

\subsection{A regression that makes destruction precede submission}

The permanent regression test is deliberately stricter than ``a short-lived target does not
crash.'' In \texttt{modules/renderers/sdl-gpu/examples/\allowbreak{}sdlgpu\_rendertarget\_lifetime\_test.cpp}, one \(8\times8\)
\cnaclass{RenderTarget2D} is created as a local variable inside a single \texttt{Draw()} call,
cleared red, then sampled by a queued \cnaclass{SpriteBatch} draw into a still-live
\(16\times16\) destination target that was first cleared blue. The local variable reaches the
end of its scope before \texttt{Present()} invokes \texttt{EnsureFrameRendered()}, so the
sampled draw is still only a command waiting to be submitted:

\begin{lstlisting}
void DrawThroughShortLivedRenderTarget(GraphicsDevice& dev) {
    RenderTarget2D localRt(dev, 8, 8, false, SurfaceFormat::Color,
                           DepthFormat::None, 0, RenderTargetUsage::DiscardContents);
    dev.SetRenderTarget(&localRt);
    dev.Clear(Color::Red);
    dev.SetRenderTarget(nullptr);

    dev.SetRenderTarget(rtDest_.get());
    dev.Clear(Color::Blue);
    sb_->Begin();
    sb_->Draw(localRt, Rectangle(0, 0, 16, 16),
              Rectangle(0, 0, 8, 8), Color::White);
    sb_->End();
    dev.SetRenderTarget(nullptr);
} // localRt is destroyed here, before this frame is submitted
\end{lstlisting}

On the first frame the test reads the centre pixel of \texttt{rtDest\_} back and requires exact
red (within its small channel tolerance), proving that the queued sampling operation really
ran; a no-crash-only assertion could otherwise pass after silently omitting the draw. It then
performs the same create--clear--sample--destroy pattern for 120 frames, which checks that the
deferred-release queue is drained after each submission rather than merely retaining every
texture until shutdown. The implementation has two necessary lifetime stages: queued targets
and draw targets retain a \texttt{shared\_ptr} to their internal state until frame rendering has
finished, then that state's destructor places its raw color, depth, and MSAA handles in
\texttt{pendingTextureReleases\_}. After a successful command-buffer submission, the renderer
releases each queued handle through \texttt{SDL\_ReleaseGPUTexture()}. SDL may still fence the
physical memory internally; CNA's narrower rule is that a recording cannot release a handle
which an unsubmitted command still names.

\section{MRT is a real shader capability, not merely a multi-clear convenience}

\cnaclass{GraphicsDevice::SetRenderTargets()} makes the first
\cnaclass{RenderTarget2D} the primary target and records every later target as one simultaneous
SDL\_gpu color attachment. At \texttt{Present()} the primary's render pass contains
\(1+\lvert\texttt{mrtSiblings}\rvert\) \texttt{SDL\_GPUColorTargetInfo} entries; the custom
effect pipeline cache includes that attachment count, so a pipeline created for one color output
cannot accidentally be rebound for two. All attachments receive the same CNA blend state, just
as one \cnaclass{GraphicsDevice::BlendState} governs the whole draw rather than one state per
target.

That does \emph{not} mean every draw writes every attachment. CNA's stock sprite and 3D fragment
shaders declare one output, so their real, deliberately narrow MRT contract is: target zero
receives the draw while targets one onward can still be bound and cleared independently. The
second half of \texttt{sdlgpu\_mrt\_test.cpp} proves the broader custom-effect path by binding
two fresh targets and issuing exactly one \cnaclass{SpriteBatch} draw with a \texttt{\#version
450} fragment shader containing \texttt{layout(location=0) out vec4 outColorA} and
\texttt{layout(location=1) out vec4 outColorB}. Sampling a white texture with tint
\((0.2,0.4,0.8,1)\) must read \((51,102,204,255)\) from A and its shader-computed channel swap
\((102,204,51,255)\) from B. Those distinct readbacks rule out both an extra clear and a copied
single-output result: one fragment invocation really wrote two simultaneous attachments.

MRT lifetime follows the same ownership repair as ordinary deferred targets. Each immutable
\texttt{PassSegment} retains its primary and every secondary attachment as
\texttt{shared\_ptr<\allowbreak{}SdlGpuRenderTarget2DState>}; explicit clears operate on the segment rather
than a parallel wrapper-pointer list. Destroying a public secondary wrapper therefore cannot
invalidate an already-recorded bind cycle, and the render-target-lifetime regression covers the
create--bind--draw/clear--destroy--present sequence.

\section{Usage selects per-resource first-use load/store behavior}

Each 2D target and each cube face begins with independent first-use colour/depth/stencil flags.
The shared Discard bind records the FNA-style black/max-depth/zero-stencil clear into its new
bind-cycle segment; Preserve records none. Once a pass has written the resource, later cycles
load its stored aspects unless an explicit ordered clear says otherwise. MRT applies the
primary cycle's explicit colour clear to all simultaneous colour attachments while retaining
each attachment's owning state. The depth/stencil usage fixture renders distinct occlusion and
stencil-gate outcomes across unbind/rebind cycles, so the result is not inferred from colour.

Cube usage now reaches construction too. A single-sample face naturally loads its cube layer;
a multisampled face owns a persistent per-face attachment and uses
\texttt{RESOLVE\_AND\_STORE} for Preserve so the next cycle can load its samples. Discard may
use plain \texttt{RESOLVE} when no later segment in the frame needs them. The six-face MSAA/mip
and depth/stencil usage suites cover these paths. \texttt{PlatformContents} follows the shared
FNA predicate \texttt{usage != DiscardContents} and therefore preserves like Preserve; cube
and depth/stencil usage fixtures assert that policy explicitly. See
\S\ref{sec:backend-rendertarget-usage-matrix}.

\subsection{A worked example: \cnaclass{Texture3D}'s own \texttt{cycle=true} orphaned-write bug}

\cnaclass{Texture3D} support on this renderer --- sub-volume upload/readback plus mipmap generation
--- surfaced a real data-corruption bug of exactly the kind this book keeps returning to: code that compiles cleanly, throws
nothing, and produces a wrong answer silently. The relevant upload operation is
\texttt{SDL\_UploadToGPUTexture}. Its \texttt{cycle} parameter offers a real performance
trade-off SDL3's own documentation states plainly: \texttt{cycle=true} swaps the texture to a
fresh, separate underlying GPU resource on each call, avoiding a stall if the GPU is still
reading the previous resource's content. That is correct and desirable for \cnaclass{Texture2D}'s
single-full-texture-replace \texttt{SetData()} path. CNA's SDL\_GPU texture renderer calls that
path \texttt{UpdatePixels}; it deliberately retains \texttt{cycle=true}. The first
\cnaclass{Texture3D} implementation copied that choice without accounting for a 3D texture's
multiple writes.
The assumption breaks for \cnaclass{Texture3D} specifically, because its content is built up
through \emph{multiple independent} sub-volume and per-level \texttt{SetData()} calls that all
need to land on the \emph{same} underlying resource --- with \texttt{cycle=true}, an earlier
partial write (say, a sub-volume at $z=0$) silently orphaned itself onto an abandoned GPU
resource the moment a second, later write (a sub-volume at $z=1$) followed it, reading back as
zero/uninitialized rather than the real content actually uploaded.

The real regression test that caught this (\texttt{modules/renderers/sdl-gpu/examples/\allowbreak{}sdlgpu\_texture3d\_test.cpp})
is deliberately shaped to make the bug unmissable rather than merely possible to catch:

\begin{lstlisting}
Texture3D tex(dev, 4, 4, 2, false, SurfaceFormat::Color);
// A 2x2x2 sub-volume at offset (1,1,0) -- red at z=0, green at z=1 -- deliberately
// off-center within the full 4x4x2 volume, and two SEPARATE SetData() calls, one per Z slice.
tex.SetData(0, /*x=*/1, /*y=*/1, /*z=*/0, 2, 2, 1, redPixels.data(), redPixels.size());
tex.SetData(0, /*x=*/1, /*y=*/1, /*z=*/1, 2, 2, 1, greenPixels.data(), greenPixels.size());
// With cycle=true, the z=0 write above would already be silently orphaned by the time this
// GetData() runs, reading back zero instead of red.
std::vector<Color> gotRed(4), gotGreen(4);
tex.GetData(0, 1, 1, 0, 2, 2, 1, gotRed.data(), gotRed.size() * sizeof(Color));
tex.GetData(0, 1, 1, 1, 2, 2, 1, gotGreen.data(), gotGreen.size() * sizeof(Color));
\end{lstlisting}

Two design choices make this test genuinely discriminating rather than accidentally passing
either way: the sub-volume is off-center (a bug that only corrupted, say, the texture's
origin corner would still show up), and each Z slice gets its own distinct solid color rather
than reusing one color for the whole volume --- a uniform-color test cannot tell correctly
-downsampled or correctly-retained content apart from an out-of-range read that happens to
land on more of that same color, the identical discriminating-test principle
Section~\ref{sec:sdlgpu-what-remains-open}'s own generated-mipmap gap needs applied to close it.
The fix, confirmed against the real
\texttt{SdlGpuTexture3DRenderer::\allowbreak{}SetData} implementation
directly, is a one-word change: pass \texttt{cycle=false} for every \cnaclass{Texture3D}
upload, leaving \cnaclass{Texture2D}'s own \texttt{cycle=true} single-replace path untouched
and correct as-is. The identical fix (and the identical reasoning, cited by name in that
renderer class's own source comment) was applied proactively to
\cnaclass{TextureCube}'s per-face uploads too, since a cube map is built up through six
separate per-face \texttt{SetData()} calls that must all land on the same underlying resource --- caught
and fixed before shipping rather than found by reproducing the same bug a second time.

\section{Forced swapchain failure preserves the queued frame}

SDL\_gpu distinguishes two acquisition outcomes that a renderer must not conflate. A successful
\texttt{SDL\_WaitAndAcquireGPUSwapchainTexture()} call with a null texture is a documented,
non-error case such as a minimized window: CNA submits the command buffer and simply skips that
frame. A false return is a hard acquisition failure. SDL forbids cancelling the command buffer
after this call, so CNA submits it first and then throws with the captured SDL error. Crucially,
it leaves \texttt{framePending\_} true: a caller which recovers the window must be able to submit
the same queued clear and draws rather than silently lose them.

\texttt{sdlgpu\_\allowbreak swapchain\_\allowbreak recovery\_\allowbreak test.cpp} proves that latter path without pretending a real
device loss is easy to reproduce. On frame 10 it calls SDL's own
\texttt{SDL\_\allowbreak ReleaseWindow\allowbreak FromGPU\allowbreak Device()}, makes \texttt{Present()} fail, then
uses SDL's matching claim call to reclaim the identical window before calling \texttt{Present()}
again. The test
requires all five stages: nine ordinary frames, the induced throw, successful reclaim, successful
presentation of the preserved frame, and 30 further frames without exception. This is a bounded,
controlled proof of recovery after a hard acquisition failure; it does not claim to solve the
separate lazy-proxy backbuffer readback path described below.

\section{What remains open}
\label{sec:sdlgpu-what-remains-open}

Three limits are worth stating precisely rather than glossing over. First, occlusion queries are a
\emph{permanent} limitation of the underlying API, not a task gap: the vendored
\texttt{SDL\_gpu.h} has no query-pool or occlusion-query type anywhere in it (only
\texttt{SDL\_GPUFence} for CPU/GPU synchronization), so \texttt{CreateOcclusionQuery()}
correctly returns \texttt{nullptr} on this renderer, the identical posture \texttt{HEADLESS} and
\texttt{SOFTWARE} already take for their own, unrelated reasons --- worth re-opening only if a
future SDL3 release adds real query support, not before. Second, the generic \texttt{.cnj} custom-effect fixture remains a
cross-renderer portability gap: its loose GLSL ES 3.00 declarations are not accepted by this
renderer's explicit-location, uniform-block binding contract, even though the dedicated
\texttt{\#version 450} effect regression is pixel-verified. That choice is deliberately left
open pending either a shared content-dialect migration or a real shader-translation layer.
Third, hardware instancing is absent at a different layer: this otherwise 3D-capable renderer
does not override \texttt{DrawInstancedPrimitivesEx()}, so every valid public
\texttt{DrawInstancedPrimitives()} call reaches the common
\texttt{std::runtime\_error} default. SDL\_gpu itself exposes per-instance vertex input; the
gap is CNA-side pipeline/binding work, not a permanent underlying-API limitation. The public
\cnaclass{GraphicsCapability::Instancing} entry now exists, but this renderer does not override
the permissive base policy and therefore incorrectly reports it as supported. This is a
capability overclaim: callers cannot safely use the bit as a guard until SDL\_GPU either
implements the draw route or returns false explicitly
(Section~\ref{sec:drawex-instancing-matrix}).

Five former entries no longer belong in this list. Depth bias and slope-scale bias are
normalized per polygon topology, hashed into every relevant pipeline key, and applied through
SDL's pipeline-static rasterizer fields; a dedicated pixel/validation test covers them.
\texttt{MaxAnisotropy} is part of the complete sampler key and native descriptor. Secondary MRT
attachments are retained as \texttt{shared\_ptr}s in the immutable pass segment, closing the raw
pointer lifetime hazard. Finally, ordinary \cnaclass{Texture2D} allocates its declared native
mip count and overrides \texttt{UpdatePixelsLevel}; the native-allocation/upload regression
proves that a nonzero level reaches the GPU resource. None of these closures adds automatic
downsampling to plain \cnaclass{Texture2D}: its upper levels remain explicitly authored, as in
the XNA/FNA contract.

Backbuffer readback is also closed without charging games that never use it. SDL's swapchain
texture remains permanently write-only, but the first \texttt{GetBackBufferData()} lazily enables
a self-owned, swapchain-sized proxy. The pending frame renders into that proxy, one 1:1 blit
presents it, and \texttt{ReadBackbuffer()} downloads the requested logical rectangle through a
transfer buffer and fence, converting BGRA to public RGBA when needed. Before the first read the
proxy and blit do not exist. First-read, whole-logical-dimension, and fully-written-range tests
cover the observable contract.

Beyond these three remaining limits, Windows (D3D12 driver) and macOS/iOS (Metal driver)
remain code paths only, not validation claims, until run on real (or Wine/DXVK-class) hardware
--- the same phased-claim discipline this book already applies to \texttt{D3D11} /
\texttt{D3D12} (Chapter~\ref{ch:direct3d-backends}) and \texttt{WebGPU}
(Chapter~\ref{ch:webgpu-backend}).

\section{Clear state is per segment and issue-ordered}

This renderer carries independent colour, depth, and stencil requests on each backbuffer,
2D-target/MRT, or cube-face segment. Consecutive clears before any draw coalesce safely into the
segment's load action, aspect by aspect. Once a draw has made ordering observable, the next
clear opens another segment over the same destination; the following native pass loads
unselected aspects, clears selected ones, and then replays later draws. MSAA segments use
\texttt{RESOLVE\_AND\_STORE} whenever a later segment must load their samples.

The renderer-neutral clear-options and ordered-clear fixtures distinguish isolated aspects,
draw--clear versus clear--draw, multiple clears, and backbuffer/2D/cube/MRT destinations without
an intermediate flush. Pass-boundary and backbuffer-order suites separately prove that target
rebinding and target--backbuffer interleaving remain in public order. Deferred submission is
therefore an implementation strategy, not a remaining clear-order exception; the comparative
matrix is \S\ref{sec:backend-clear-matrix}.

\section{Verification methodology}

The executable claims above have dedicated tests rather than resting on code reads alone.
Counting the actual module-local registrations directly
(\texttt{modules/renderers/sdl-gpu/\allowbreak{}examples/CMakeLists.txt}) rather than trusting an old plan
snapshot shows \textbf{83} \texttt{SdlGpu\_*} CTest registrations today, spanning the 2D and 3D
smoke tests, every stock effect (including \cnaclass{EnvironmentMapEffect},
\cnaclass{SkinnedEffect}, and the \texttt{CNAEXT} \cnaclass{PbrEffect} this book's WebGPU chapter
already introduced), custom \cnaclass{ShaderEffect}, every render-target variant (2D, cube,
multiple render targets, MSAA), both extra texture types, sampler/render state, draw order, and
swapchain recovery, plus the render-target-lifetime regression above --- run, like every other renderer's own CTest suite, against this project's
own \texttt{cna\_register\_renderer\_test()} convention, gracefully skipping with a clear reason
rather than hard-failing when no real \texttt{DISPLAY} / \texttt{WAYLAND\_DISPLAY} is present, the
identical convention \texttt{D3D11} / \texttt{D3D12} / \texttt{WebGPU} already follow. All verification
to date is against this project's own real dev-machine GPU via \texttt{SDL\_gpu}'s Vulkan
driver, the only driver this project's own prebuilt SDL3 packages compile in on Linux --- the
same single-driver-verified, other-drivers-unverified posture this chapter's own ``What remains
open'' section already names honestly rather than implying broader coverage than actually
exists.
